\documentclass{article}

\usepackage{listings}
\usepackage{fontspec}
\usepackage{polyglossia}
\setmainlanguage{greek}
\setotherlanguage{english}
\setmainfont{FreeSerif}
\setmonofont{FreeMono}
\usepackage[a4paper,top=2cm,bottom=2cm,left=3cm,right=3cm,marginparwidth=1.75cm]{geometry}

\usepackage{amsmath}
\usepackage{amssymb}
\usepackage{graphicx}
\usepackage[colorlinks=true, allcolors=blue]{hyperref}
\usepackage{float}
\usepackage{xcolor}

\definecolor{dkgreen}{rgb}{0,0.6,0}
\definecolor{mauve}{rgb}{0.58,0,0.82}

\title{Εργασία \#2: Μετασχηματισμοί και Προβολές}
\author{Βογιατζής Χαρίσιος ΑΕΜ:9192}

\lstset{frame=tb,
  language=Python,
  aboveskip=3mm,
  belowskip=3mm,
  showstringspaces=false,
  columns=flexible,
  basicstyle={\small\ttfamily},
  numbers=none,
  numberstyle=\tiny\color{gray},
  keywordstyle=\color{blue},
  commentstyle=\color{dkgreen},
  stringstyle=\color{mauve},
  breaklines=true,
  breakatwhitespace=true,
  tabsize=3
}

\begin{document}
\maketitle

\begin{center}
\href{https://github.com/charisvt/cg-auth-hw2-2026}{github.com/charisvt/cg-auth-hw2-2026}
\end{center}

\begin{abstract}
Σκοπός της εργασίας είναι η υλοποίηση του pipeline προβολής τρισδιάστατων αντικειμένων σε
δισδιάστατη εικόνα, καθώς και η παραγωγή δύο animations: ένα με στατική κάμερα και
περιστρεφόμενο αντικείμενο, και ένα με κάμερα που τροχιάζει γύρω από ένα στατικό αντικείμενο.
\end{abstract}

% -----------------------------------------------------------------------
\section{Δομή κώδικα και τρόπος κλήσης}

Ο κώδικας οργανώνεται στα παρακάτω αρχεία:
\begin{itemize}
  \item \texttt{transformations.py} --- κλάση \texttt{Trafo} και συναρτήσεις
        \texttt{lookat}, \texttt{perspective\_project}, \texttt{rasterize},
        \texttt{render\_object}.
  \item \texttt{renderer.py} --- βοηθητικές συναρτήσεις Gouraud shading από
        την εργασία \#1 (\texttt{vector\_interp}, \texttt{g\_shading}).
  \item \texttt{demo1.py} --- παράγει 25 frames με στατική κάμερα και
        περιστρεφόμενο αντικείμενο.
  \item \texttt{demo2.py} --- παράγει 25 frames με κινούμενη κάμερα γύρω από
        το στατικό αντικείμενο.
\end{itemize}

Εκτέλεση των demos (απαιτείται Python $\geq$ 3.12, τα frames αποθηκεύονται στους
φακέλους \texttt{frames\_demo1/} και \texttt{frames\_demo2/} αντίστοιχα):
\begin{lstlisting}
python demo1.py
python demo2.py
\end{lstlisting}

% -----------------------------------------------------------------------
\section{Κλάση \texttt{Trafo}}

Η κλάση \texttt{Trafo} υλοποιεί έναν affine μετασχηματισμό (περιστροφή + μετατόπιση)
μέσω ενός $3\times3$ πίνακα περιστροφής \texttt{rot\_mat} και ενός διανύσματος
μετατόπισης \texttt{t\_vec}.

\subsection{\texttt{translate}}
Προσθέτει το δοθέν διάνυσμα στο \texttt{t\_vec}:
\begin{lstlisting}
def translate(self, t_vec: np.ndarray) -> None:
    self.t_vec = self.t_vec + t_vec
\end{lstlisting}

\subsection{\texttt{rotate}}
Υπολογίζει τον πίνακα περιστροφής με τον τύπο Rodrigues (σχ. 5.45 σημειώσεων) για
περιστροφή γωνίας $a$ γύρω από μοναδιαίο άξονα $\mathbf{u} = [u_x, u_y, u_z]^\top$:
\[
\mathbf{R} = (1-\cos a)
\begin{bmatrix} u_x^2 & u_x u_y & u_x u_z \\ u_y u_x & u_y^2 & u_y u_z \\ u_z u_x & u_z u_y & u_z^2 \end{bmatrix}
+ \cos a\, \mathbf{I}_{3\times3}
+ \sin a
\begin{bmatrix} 0 & -u_z & u_y \\ u_z & 0 & -u_x \\ -u_y & u_x & 0 \end{bmatrix}
\]
Επειδή ο άξονας διέρχεται από το \texttt{center}, το \texttt{t\_vec}
περιστρέφεται ως προς αυτό πριν ενημερωθεί:
\begin{lstlisting}
self.t_vec   = np.dot(R_new, self.t_vec - center) + center
self.rot_mat = np.dot(R_new, self.rot_mat)
\end{lstlisting}

\subsection{\texttt{xform\_pnts}}
Εφαρμόζει τον affine μετασχηματισμό στα σημεία εισόδου ($3\times N$):
\begin{lstlisting}
def xform_pnts(self, pts: np.ndarray) -> np.ndarray:
    return np.dot(self.rot_mat, pts) + self.t_vec[:, np.newaxis]
\end{lstlisting}

\subsection{\texttt{sys2sys}}
Το \texttt{sys\_src} $\in\mathbb{R}^{3\times3}$ είναι ο πίνακας περιστροφής
που ορίζει το σύστημα συντεταγμένων πηγής (αντίστοιχο της περίπτωσης 2 της
§5.4.3 σημειώσεων), ενώ \texttt{pts} είναι τα σημεία εκφρασμένα σε αυτό.
Η μετατροπή γίνεται σε δύο βήματα: πηγή $\to$ WCS, στη συνέχεια WCS $\to$
σύστημα προορισμού ($R^{-1}=R^\top$ για ορθογώνιο $R$):
\begin{lstlisting}
pts_world = np.dot(sys_src, pts)
return np.dot(self.rot_mat.T, pts_world - self.t_vec[:, np.newaxis])
\end{lstlisting}

% -----------------------------------------------------------------------
\section{Συνάρτηση \texttt{lookat}}

Υπολογίζει τον πίνακα $R \in \mathbb{R}^{3\times3}$ και το διάνυσμα
$\mathbf{t} \in \mathbb{R}^3$ που μετασχηματίζουν σημεία από το WCS στο σύστημα
της κάμερας (§6.3.1--6.3.2 σημειώσεων). Με $C = \mathtt{eye}$, $K = \mathtt{target}$,
$\mathbf{u} = \mathtt{up}$:

\begin{enumerate}
  \item $\hat{z}_c = \dfrac{\overrightarrow{CK}}{|\overrightarrow{CK}|}$ \quad (άξονας φακού, εξ. 6.6)
  \item $\mathbf{t} = \mathbf{u} - \langle\mathbf{u},\hat{z}_c\rangle\hat{z}_c,\quad
        \hat{y}_c = \dfrac{\mathbf{t}}{|\mathbf{t}|}$ \quad (διόρθωση up, εξ. 6.7)
  \item $\hat{x}_c = \hat{y}_c \times \hat{z}_c$ \quad (εξ. 6.8)
  \item $R = \begin{bmatrix}\hat{x}_c^\top\\\hat{y}_c^\top\\\hat{z}_c^\top\end{bmatrix}, \quad
        \mathbf{t} = -R\,C$ \quad (εξ. 6.11, 6.13)
\end{enumerate}

Ο $R$ έχει τα διανύσματα βάσης της κάμερας ως γραμμές (ισοδύναμο με τον $R^\top$
των σημειώσεων όπου τα διανύσματα βάσης είναι στήλες). Τα σημεία μπροστά από
την κάμερα έχουν $Z_c > 0$.

\begin{lstlisting}
R = np.array([right, up_cam, forward])
t = -np.dot(R, eye)
return R, t
\end{lstlisting}

% -----------------------------------------------------------------------
\section{Συνάρτηση \texttt{perspective\_project}}

Μετασχηματίζει τα σημεία στο σύστημα της κάμερας και εφαρμόζει την
προοπτική προβολή pinhole (§6.2.1, εξ. 6.2 σημειώσεων):
\[
p_c = R\,p + \mathbf{t}, \qquad
x_q = w\,\frac{x_p}{z_p}, \quad y_q = w\,\frac{y_p}{z_p}
\]
όπου $w$ είναι η απόσταση κέντρου προβολής -- πετάσματος. Επιστρέφει τις
2D προβολές και το βάθος $z_p$ για την ταξινόμηση κατά βάθος.

\begin{lstlisting}
p_cam  = np.dot(R, pts) + t[:, np.newaxis]
depth  = p_cam[2]
pts_2d = focal * p_cam[:2] / p_cam[2:3]
return pts_2d, depth
\end{lstlisting}

% -----------------------------------------------------------------------
\section{Συνάρτηση \texttt{rasterize}}

Μετατρέπει τις 2D συντεταγμένες από το επίπεδο πετάσματος
($[-w/2, w/2]\times[-h/2, h/2]$) σε pixel εικόνας. Ο άξονας της κάμερας
αντιστοιχεί στο κέντρο της εικόνας, ενώ ο $y$ αντιστρέφεται επειδή η
αρίθμηση ξεκινά από την πάνω αριστερή γωνία:
\[
x_{px} = \left(\frac{x}{w} + 0.5\right)w_{res}, \qquad
y_{px} = \left(0.5 - \frac{y}{h}\right)h_{res}
\]

\begin{lstlisting}
x_pix = (pts_2d[0] / plane_w + 0.5) * res_w
y_pix = (0.5 - pts_2d[1] / plane_h) * res_h
return np.round(np.stack([x_pix, y_pix], axis=1)).astype(int)
\end{lstlisting}

% -----------------------------------------------------------------------
\section{Συνάρτηση \texttt{render\_object}}

Υλοποιεί το πλήρες pipeline απεικόνισης: \texttt{lookat} $\to$
\texttt{perspective\_project} $\to$ \texttt{rasterize} $\to$ \texttt{g\_shading}.
Τα τρίγωνα ταξινομούνται από πίσω προς τα μπρος βάσει του μέσου βάθους
των κορυφών τους (Painter's Algorithm), και για κάθε τρίγωνο καλείται η
\texttt{g\_shading} από την εργασία \#1.

\begin{lstlisting}
img = np.ones((res_h, res_w, 3), dtype=np.float64)
avg_depth = np.mean(depth[t_pos_idx], axis=1)
order = np.argsort(avg_depth)[::-1]

for idx in order:
    face = t_pos_idx[idx]
    img = g_shading(img, vertices_px[face], v_clr[face])

return img
\end{lstlisting}

% -----------------------------------------------------------------------
\section{Demos}

Και τα δύο demos φορτώνουν τις παραμέτρους από το αρχείο \texttt{data.npy}
(κάμερα, αντικείμενο, παράμετροι animation). Τα χρώματα των κορυφών
(\texttt{v\_clr}) προκύπτουν με δειγματοληψία της εικόνας texture
\texttt{loony-repeat.png} στις UV συντεταγμένες κάθε κορυφής.
Κάθε demo παράγει 25 frames ($25\,\text{fps} \times 1\,\text{sec}$).

\subsection{Demo 1 --- Στατική κάμερα, περιστρεφόμενο αντικείμενο}
Για κάθε frame $i$ δημιουργείται νέο \texttt{Trafo} με γωνία
$\theta = 2\pi i/25$ γύρω από τον άξονα $Y$ (κέντρο αρχή αξόνων).
Η κάμερα παραμένει σταθερή:
\begin{lstlisting}
trafo = Trafo()
trafo.rotate(np.array([0.0, 1.0, 0.0]), angle, np.array([0.0, 0.0, 0.0]))
v_pos_rot = trafo.xform_pnts(v_pos)
img = render_object(v_pos_rot, v_clr, t_pos_idx, ...)
\end{lstlisting}

\subsection{Demo 2 --- Περιστρεφόμενη κάμερα, στατικό αντικείμενο}
Η κάμερα κινείται σε κύκλο ακτίνας \texttt{k\_cam\_radius} στο επίπεδο
$XZ$, κοιτώντας πάντα το αντικείμενο. Η δεξιόστροφη περιστροφή (ως προς
$+Y$) εφαρμόζεται με $R_y(-\theta)$:
\[
x' = x_0\cos\theta - z_0\sin\theta, \qquad z' = x_0\sin\theta + z_0\cos\theta
\]
\begin{lstlisting}
cam_eye = np.array([x0*ct - z0*st, y0, x0*st + z0*ct])
img = render_object(v_pos, v_clr, t_pos_idx, ..., eye=cam_eye, ...)
\end{lstlisting}

% -----------------------------------------------------------------------
\section{Αποτελέσματα}

\begin{figure}[H]
\centering
\includegraphics[width=0.45\textwidth]{frames_demo1/frame_000.png}
\includegraphics[width=0.45\textwidth]{frames_demo1/frame_006.png}
\caption{Demo 1 --- frames 0 και 6 (περιστροφή αντικειμένου, στατική κάμερα).}
\end{figure}

\begin{figure}[H]
\centering
\includegraphics[width=0.45\textwidth]{frames_demo2/frame_000.png}
\includegraphics[width=0.45\textwidth]{frames_demo2/frame_006.png}
\caption{Demo 2 --- frames 0 και 6 (περιστροφή κάμερας, στατικό αντικείμενο).}
\end{figure}

% -----------------------------------------------------------------------
\section{Δημιουργία video από τα frames}

Τα αποθηκευμένα frames μπορούν να συνδυαστούν σε βίντεο με \texttt{ffmpeg}:
\begin{lstlisting}
ffmpeg -r 25 -i frames_demo1/frame_%03d.png -vcodec libx264 demo1.mp4
ffmpeg -r 25 -i frames_demo2/frame_%03d.png -vcodec libx264 demo2.mp4
\end{lstlisting}

% -----------------------------------------------------------------------
\section{Κώδικας}
Μπορείτε να βρείτε τον κώδικα και τα αποτελέσματα στο \href{https://github.com/charisvt/cg-auth-hw2-2026}{Github}.

\end{document}
